\documentclass[12pt,a4paper]{article}

% ── Encodage & langue ─────────────────────────────────────────────────────────
\usepackage[utf8]{inputenc}
\usepackage[T1]{fontenc}
\usepackage[french]{babel}


% ── Mise en page ──────────────────────────────────────────────────────────────
\usepackage[top=2.5cm, bottom=2.5cm, left=2.8cm, right=2.8cm]{geometry}
\usepackage{setspace}
\onehalfspacing
\usepackage{fancyhdr}
\pagestyle{fancy}
\fancyhf{}
\fancyhead[L]{\small\textit{Optimisation de géométrie de mur anti-radar}}
\fancyhead[R]{\small\textit{FDTD 2D + Algorithmes d'apprentissage}}
\fancyfoot[C]{\thepage}
\renewcommand{\headrulewidth}{0.4pt}

% ── Mathématiques ─────────────────────────────────────────────────────────────
\usepackage{amsmath, amssymb, amsthm}
\usepackage{mathtools}
\usepackage{empheq}
\usepackage{bm}         % gras pour vecteurs

% ── Figures & couleurs ────────────────────────────────────────────────────────
\usepackage{graphicx}
\usepackage{xcolor}
\usepackage{tikz}
\usetikzlibrary{arrows.meta, shapes.geometric, positioning, decorations.pathmorphing,
                patterns, calc, matrix, fit}
\usepackage{pgfplots}
\pgfplotsset{compat=1.18}

% ── Tableaux ──────────────────────────────────────────────────────────────────
\usepackage{booktabs}
\usepackage{array}
\usepackage{multirow}

% ── Algorithmes ───────────────────────────────────────────────────────────────
\usepackage{algorithm}
\usepackage{algpseudocode}
\algnewcommand\algorithmicinput{\textbf{Entrée :}}
\algnewcommand\algorithmicoutput{\textbf{Sortie :}}
\algnewcommand\Input{\item[\algorithmicinput]}
\algnewcommand\Output{\item[\algorithmicoutput]}

% ── Code source ───────────────────────────────────────────────────────────────
\usepackage{listings}
\definecolor{codeblue}{RGB}{0,0,150}
\definecolor{codegray}{RGB}{128,128,128}
\definecolor{codegreen}{RGB}{0,128,0}
\definecolor{backcolour}{RGB}{245,245,245}
\lstdefinestyle{pythonstyle}{
    backgroundcolor=\color{backcolour},
    commentstyle=\color{codegreen}\itshape,
    keywordstyle=\color{codeblue}\bfseries,
    numberstyle=\tiny\color{codegray},
    stringstyle=\color{red!70!black},
    basicstyle=\ttfamily\small,
    breakatwhitespace=false, breaklines=true,
    captionpos=b, keepspaces=true,
    numbers=left, numbersep=5pt,
    showspaces=false, showstringspaces=false,
    showtabs=false, tabsize=2,
    language=Python,
    frame=single, rulecolor=\color{gray!50}
}
\lstset{style=pythonstyle}

% ── Références ────────────────────────────────────────────────────────────────
\usepackage[colorlinks=true, linkcolor=blue!70!black,
            citecolor=green!50!black, urlcolor=blue!70!black]{hyperref}
\usepackage{cleveref}

% ── Boîtes ────────────────────────────────────────────────────────────────────
\usepackage{tcolorbox}
\tcbuselibrary{theorems, skins, breakable}
\newtcolorbox{defbox}[1]{
    colback=blue!5, colframe=blue!40!black,
    fonttitle=\bfseries, title=#1, breakable
}
\newtcolorbox{resultbox}[1]{
    colback=green!5, colframe=green!40!black,
    fonttitle=\bfseries, title=#1, breakable
}
\newtcolorbox{remarque}[1]{
    colback=orange!5, colframe=orange!60!black,
    fonttitle=\bfseries, title=#1, breakable
}

% ── Notation ──────────────────────────────────────────────────────────────────
\newcommand{\vE}{\bm{E}}
\newcommand{\vH}{\bm{H}}
\newcommand{\vB}{\bm{B}}
\newcommand{\vD}{\bm{D}}
\newcommand{\vJ}{\bm{J}}
\newcommand{\vn}{\hat{\bm{n}}}
\newcommand{\Ez}{{E_z}}
\newcommand{\Hx}{{H_x}}
\newcommand{\Hy}{{H_y}}
\newcommand{\eps}{\varepsilon}
\newcommand{\Sc}{S_c}
\newcommand{\dx}{\Delta x}
\newcommand{\dt}{\Delta t}
\newcommand{\half}{\tfrac{1}{2}}

% ── Numérotation équations par section ────────────────────────────────────────
\numberwithin{equation}{section}
\numberwithin{figure}{section}
\numberwithin{table}{section}

% ══════════════════════════════════════════════════════════════════════════════
\begin{document}

% ── Page de titre ─────────────────────────────────────────────────────────────
\begin{titlepage}
\centering
\vspace*{1cm}

{\Large \textsc{CY Cergy Paris University}}\\[0.5cm]
{\large Annexe simulation}\\[2cm]

\hrule height 1.5pt \vspace{0.4cm}
{\LARGE \bfseries Optimisation de géométrie de mur anti-radar}\\[0.3cm]
{\Large \bfseries par simulation électromagnétique FDTD 2D}\\[0.2cm]
{\Large \bfseries et algorithmes d'apprentissage automatique}\\[0.3cm]
\hrule height 1.5pt

\vspace{1.5cm}

\begin{tabular}{ll}
\textbf{Auteur} & Alexis Leventoux \\[0.2cm]
\textbf{Niveau} & L3 --- Physique\\[0.2cm]
\textbf{Date} & Mars 2026 \\
\end{tabular}

\vspace{1.5cm}

\begin{abstract}
\noindent
Ce rapport présente la conception, la modélisation physique et l'implémentation
numérique d'un système d'optimisation de surface anti-radar.
L'objectif est de minimiser la \emph{Section Efficace Radar} (SER, ou RCS en anglais)
d'un mur conducteur parfait (PEC) par action sur sa géométrie.

La modélisation repose sur la méthode \emph{Finite-Difference Time-Domain}
(FDTD) en polarisation TMz, qui résout les équations de Maxwell sur une grille
de Yee 2D. La source est injectée par la technique
\emph{Total-Field / Scattered-Field} (TFSF), les frontières sont absorbées par
une couche \emph{Convolutional PML} (CPML), et la SER est calculée via une
transformation champ proche --- champ lointain (NTFF).

L'optimisation de la forme du mur est assurée par un \emph{Algorithme Génétique} (GA) avec croisement SBX et
mutation polynomiale adaptive (règle du 1/5 de Rechenberg). Des réductions de SER supérieures
à $\mathbf{90\,\%}$ sont obtenues par rapport à un mur plat de référence.
\end{abstract}

\vfill
\end{titlepage}

% ── Table des matières ────────────────────────────────────────────────────────
\tableofcontents
\newpage

% ══════════════════════════════════════════════════════════════════════════════
\section{Introduction et contexte physique}
% ══════════════════════════════════════════════════════════════════════════════

\subsection{Motivation : la furtivité radar}

La furtivité radar consiste à minimiser la quantité d'énergie électromagnétique
renvoyée vers la source après illumination d'un objet.
Deux mécanismes complémentaires permettent d'atteindre cet objectif~:
\begin{itemize}
    \item \textbf{Composition matérielle~:} utilisation de matériaux
    absorbants radar (RAM, \textit{Radar Absorbing Materials}) qui dissipent
    l'énergie incidente par pertes diélectriques ou magnétiques.
    \item \textbf{Géométrie de la surface~:} mise en forme de la surface
    réfléchissante pour rediriger les ondes diffusées dans des directions
    non dangereuses (loin de l'émetteur radar).
\end{itemize}

La présentation expérimentale a montré, en bande X ($f_0 = 10\,\text{GHz}$, $\lambda \approx 3\,\text{cm}$),
qu'une surface aluminium arrondie réduit déjà le signal rétrodiffusé de 39\,\%
par rapport à une surface plane. Ce projet explore jusqu'où une géométrie
\emph{optimisée par algorithme} peut aller.

\subsection{Grandeurs caractéristiques du problème}

Le tableau~\ref{tab:params} résume les paramètres physiques du domaine simulé.

\begin{table}[h!]
\centering
\caption{Paramètres physiques de la simulation de référence.}
\label{tab:params}
\begin{tabular}{lll}
\toprule
\textbf{Grandeur} & \textbf{Valeur} & \textbf{Origine} \\
\midrule
Fréquence centrale & $f_0 = 10\,\text{GHz}$ & Bande X (radar classique) \\
Longueur d'onde & $\lambda = c/f_0 \approx 30\,\text{mm}$ & $c = 299\,792\,458\,\text{m/s}$ \\
Impédance du vide & $\eta_0 = \sqrt{\mu_0/\eps_0} \approx 377\,\Omega$ & \\
Résolution spatiale & $\dx = \lambda / 15 \approx 2\,\text{mm}$ & 15 pts$/\lambda$ \\
Pas temporel & $\dt = 0{,}5\,\dx/c \approx 3{,}3\,\text{ps}$ & Critère CFL \\
Domaine physique & $150\dx \times 150\dx = 30 \times 30\,\text{cm}$ & Grille numérique \\
\bottomrule
\end{tabular}
\end{table}

% ══════════════════════════════════════════════════════════════════════════════
\section{Modélisation électromagnétique --- Équations de Maxwell}
% ══════════════════════════════════════════════════════════════════════════════

\subsection{Équations de Maxwell en milieu linéaire}

On part des équations de Maxwell en milieu linéaire, isotrope,
non magnétique ($\mu_r = 1$) et sans charges libres~:

\begin{align}
    \nabla \times \vE &= -\mu_0 \frac{\partial \vH}{\partial t}
    \label{eq:rot_E} \\
    \nabla \times \vH &= \eps_0 \frac{\partial \vE}{\partial t} + \vJ_s
    \label{eq:rot_H}
\end{align}

où $\vJ_s$ est la densité de courant source (ici nulle sauf à la frontière TFSF).

\subsection{Réduction en polarisation TMz}

En 2D, on suppose invariance selon $z$ ($\partial/\partial z = 0$).
La polarisation \emph{Transverse Magnétique z} (TMz) suppose que les seules
composantes non nulles sont~:
\begin{equation}
    \Ez(x,y,t), \quad \Hx(x,y,t), \quad \Hy(x,y,t)
    \label{eq:TMz_components}
\end{equation}

Le rotationnel de \eqref{eq:rot_E} et \eqref{eq:rot_H} donne alors le
système couplé TMz~:

\begin{empheq}[box=\fbox]{align}
    \frac{\partial \Hx}{\partial t} &= -\frac{1}{\mu_0}
        \frac{\partial \Ez}{\partial y}
    \label{eq:dHx_dt} \\
    \frac{\partial \Hy}{\partial t} &= +\frac{1}{\mu_0}
        \frac{\partial \Ez}{\partial x}
    \label{eq:dHy_dt} \\
    \frac{\partial \Ez}{\partial t} &= \frac{1}{\eps_0}
        \left( \frac{\partial \Hy}{\partial x} - \frac{\partial \Hx}{\partial y} \right)
    \label{eq:dEz_dt}
\end{empheq}

Ce système est \emph{hyperbolique} du premier ordre et donc peut directement
être discrétisé par différences finies.

\subsection{Conducteur parfait électrique (PEC)}

La condition aux limites sur un conducteur parfait (PEC) impose que la
composante tangentielle du champ électrique s'annule à la surface~:
\begin{equation}
    \vn \times \vE\big|_{\partial \Omega_{\text{PEC}}} = \bm{0}
    \implies
    \Ez\big|_{\text{PEC}} = 0
    \label{eq:PEC_BC}
\end{equation}
Cette condition est appliquée en forçant $\Ez = 0$ dans toutes les cellules
appartenant au mur PEC à chaque pas de temps.

% ══════════════════════════════════════════════════════════════════════════════
\section{Méthode FDTD --- Discrétisation de Yee}
% ══════════════════════════════════════════════════════════════════════════════

\subsection{Grille de Yee}

La grille de Yee est une grille décalée en espace et en temps
qui garantit~: (i) l'ordre 2 en espace et en temps ; (ii) la vérification
naturelle de $\nabla \cdot \vB = 0$ ; (iii) l'absence de modes parasites.

Les composantes sont positionnées comme suit~:
\begin{itemize}
    \item ${E_z}^n(i,j)$~: nœuds entiers en espace, pas entier en temps.
    \item ${H_x}^{n+\half}(i, j+\half)$~: décalé de $\half\dx$ en $y$,
    de $\half\dt$ en $t$.
    \item ${H_y}^{n+\half}(i+\half, j)$~: décalé de $\half\dx$ en $x$.
\end{itemize}

\begin{figure}[h!]
\centering
\begin{tikzpicture}[scale=1.2, >=stealth]

    % --- DÉFINITION DES STYLES ---
    % L'utilisation de styles permet de garder un code propre et de garantir
    % que toutes les formes sont parfaitement centrées sur leurs coordonnées.
    \tikzset{
        ez/.style={circle, draw=blue!80!black, fill=blue!40, thick, minimum size=6pt, inner sep=0pt},
        hx/.style={rectangle, draw=red!80!black, fill=red!40, thick, minimum size=6.5pt, inner sep=0pt},
        hy/.style={regular polygon, regular polygon sides=3, shape border rotate=-90, draw=green!60!black, fill=green!50, thick, minimum size=8pt, inner sep=0pt}
    }

    % --- ZONE DE SURBRILLANCE (CELLULE DE YEE) ---
    % Placé avant la grille pour ne pas cacher les lignes
    \fill[yellow!15, rounded corners=3pt] (1.9,1.9) rectangle (4.1,4.1);
    \draw[orange, thick, densely dotted, rounded corners=3pt] (1.9,1.9) rectangle (4.1,4.1);
    \node[text=orange!90!black, anchor=south] at (3, 4.1) {\footnotesize Cellule de Yee $(i,j)$};

    % --- GRILLES ---
    % Grille étendue en pointillés pour suggérer la continuité de l'espace
    \draw[gray!30, thick, dashed] (-0.5,-0.5) grid[step=2] (4.5,4.5);
    % Grille principale
    \draw[gray!60, thick] (0,0) grid[step=2] (4,4);

    % --- PLACEMENT DES NOEUDS ---
    % Ez aux noeuds entiers
    \foreach \x in {0,2,4} {
        \foreach \y in {0,2,4} {
            \node[ez] at (\x,\y) {};
        }
    }

    % Hx aux demi-noeuds y (x entier, y demi-entier)
    \foreach \x in {0,2,4} {
        \foreach \y in {1,3} {
            \node[hx] at (\x,\y) {};
        }
    }

    % Hy aux demi-noeuds x (x demi-entier, y entier)
    \foreach \x in {1,3} {
        \foreach \y in {0,2,4} {
            \node[hy] at (\x,\y) {};
        }
    }

    % --- LÉGENDE ---
    \begin{scope}[shift={(5.5, 2)}]
        % Cadre de la légende
        \draw[gray!40, fill=black!2, rounded corners=4pt] (-0.6, -1.4) rectangle (3.2, 1.4);
        
        % Composantes de la légende
        \node[ez, label=right:{\small $\Ez(i,j)$}]       at (0, 0.8) {};
        \node[hx, label=right:{\small $\Hx(i,j+\half)$}] at (0, 0) {};
        \node[hy, label=right:{\small $\Hy(i+\half,j)$}] at (0, -0.8) {};
    \end{scope}

    % --- ANNOTATIONS DE DISTANCES ---
    % Utilisation de |<->| pour des lignes de cotes professionnelles
    \draw[|<->|, thick] (0,-0.6) -- (2,-0.6) node[midway,below]{$\dx$};
    \draw[|<->|, thick] (-0.6,0) -- (-0.6,2) node[midway,left]{$\dx$};

\end{tikzpicture}
\caption{Grille de Yee en polarisation TMz. Les composantes $\Ez$, $\Hx$, $\Hy$
sont décalées d'une demi-cellule pour assurer la consistance du rotationnel discret.}
\label{fig:yee_grid}
\end{figure}

\subsection{Schéma leapfrog --- équations de mise à jour}

On discrétise \eqref{eq:dHx_dt}--\eqref{eq:dEz_dt} par un schéma
\emph{leapfrog} centré d'ordre 2 en espace et en temps~:

\begin{align}
    {H_x}^{n+\half}_{i,j+\half} &= {H_x}^{n-\half}_{i,j+\half}
        - \frac{\dt}{\mu_0\dx}
        \left( {E_z}^n_{i,j+1} - {E_z}^n_{i,j} \right)
    \label{eq:update_Hx} \\
    {H_y}^{n+\half}_{i+\half,j} &= {H_y}^{n-\half}_{i+\half,j}
        + \frac{\dt}{\mu_0\dx}
        \left( {E_z}^n_{i+1,j} - {E_z}^n_{i,j} \right)
    \label{eq:update_Hy} \\
    {E_z}^{n+1}_{i,j} &= {E_z}^n_{i,j}
        + \frac{\dt}{\eps_0\dx}
        \left( {H_y}^{n+\half}_{i+\half,j} - {H_y}^{n+\half}_{i-\half,j}
               - {H_x}^{n+\half}_{i,j+\half} + {H_x}^{n+\half}_{i,j-\half} \right)
    \label{eq:update_Ez}
\end{align}

On définit les coefficients de mise à jour~:
\begin{equation}
    C_{Hxe} = C_{Hye} = \frac{\dt}{\mu_0 \dx}, \qquad
    C_{Ezh} = \frac{\dt}{\eps_0 \dx}
\end{equation}

\subsection{Stabilité --- Critère de Courant--Friedrichs--Lewy}

Le schéma leapfrog est explicite et conditionellement stable. En 2D,
la condition de stabilité (CFL) impose~:
\begin{equation}
    \Sc = \frac{c\,\dt}{\dx} \leq \frac{1}{\sqrt{2}} \approx 0{,}707
    \label{eq:CFL}
\end{equation}

Le projet utilise $\Sc = 0{,}5$, ce qui laisse une marge de sécurité confortable.
La dispersion numérique résiduelle vaut~:
\begin{equation}
    \tilde{c}(\theta) = c \cdot \frac{2}{\Sc}\sin^{-1}
    \!\left[\Sc \sqrt{\sin^2\!\left(\frac{k_x\dx}{2}\right)
                    + \sin^2\!\left(\frac{k_y\dx}{2}\right)}\right]
    \bigg/ \sqrt{k_x^2 + k_y^2}
\end{equation}
Elle est inférieure à 1\,\% pour 15 points par longueur d'onde.

\subsection{Source de Ricker}

La source temporelle est une wavelet de Ricker (dérivée seconde de la
Gaussienne), standard en simulation FDTD pour sa bande spectrale bien définie~:
\begin{equation}
    s(t) = \left[1 - 2\pi^2 f_p^2 (t-t_d)^2\right]
           \exp\!\left[-\pi^2 f_p^2 (t-t_d)^2\right]
    \label{eq:ricker}
\end{equation}
avec $f_p = 10\,\text{GHz}$ la fréquence centrale et
$t_d = 1/f_p$ le retard permettant un démarrage à amplitude nulle.
La fréquence maximale du spectre est $f_{\max} \approx 2{,}5\,f_p$.

% ══════════════════════════════════════════════════════════════════════════════
\section{Injection de la source --- Technique TFSF}
% ══════════════════════════════════════════════════════════════════════════════

\subsection{Principe}

La technique \emph{Total-Field / Scattered-Field} (TFSF) sépare le domaine en~:
\begin{itemize}
    \item \textbf{Zone champ total} (TF) : intérieure à une frontière virtuelle
    $\Gamma_{\text{TFSF}}$, contient $\vE_{\text{tot}} = \vE_{\text{inc}} + \vE_{\text{diff}}$.
    \item \textbf{Zone champ diffusé} (SF) : extérieure, contient uniquement
    $\vE_{\text{diff}}$.
\end{itemize}

\subsection{Corrections à la frontière TFSF}

À la frontière $\Gamma_{\text{TFSF}}$, les équations de mise à jour croisent les
deux zones et doivent être corrigées. Pour une incidence normale ($\theta = 0$),
une grille 1D auxiliaire calcule le champ incident analytique.
Pour une incidence oblique d'angle $\theta$, le champ incident s'écrit~:

\begin{align}
    E_z^{\text{inc}}(x,y,t) &= s\!\left(t - \frac{x\cos\theta + y\sin\theta}{c}\right) \\
    H_x^{\text{inc}} &= -\frac{\sin\theta}{\eta_0}\,E_z^{\text{inc}} \\
    H_y^{\text{inc}} &= +\frac{\cos\theta}{\eta_0}\,E_z^{\text{inc}}
\end{align}

Les corrections à appliquer aux bords de $\Gamma_{\text{TFSF}}$ sont~:
\begin{equation}
    H_y^{n+\half}\big|_{i_0-\half, j} \mathrel{-}=
        \frac{\dt}{\mu_0\dx}\,E_z^{\text{inc},n}\big|_{i_0, j}
\end{equation}
(et de manière symétrique pour les autres bords et composantes).

\subsection{Calcul de la SER}

On définit la \emph{Section Efficace Radar} 2D normalisée~:
\begin{equation}
    \frac{\sigma_{2D}}{\lambda} = \frac{k}{4\pi}
    \frac{\left|E_z^{\text{diff}}(\phi=\pi)\right|^2}
         {\left|\hat{E}_z^{\text{inc}}\right|^2}
    \label{eq:SER_def}
\end{equation}
où $k = 2\pi/\lambda$ est le nombre d'onde,
$\hat{E}_z^{\text{inc}}$ le spectre de la source et
$E_z^{\text{diff}}(\phi=\pi)$ le champ lointain diffusé en direction de
rétrodiffusion. Ce dernier est obtenu par transformation NTFF.

\paragraph{Référence analytique.}
Pour un mur PEC plan d'hauteur $w$, l'Optique Physique 2D (Balanis, §11.3)
donne en incidence normale~:
\begin{equation}
    \frac{\sigma_{2D}}{\lambda} = \frac{2w^2}{\lambda^2}
    \label{eq:SER_PO}
\end{equation}
Avec $w = 50\dx = 10\,\text{cm}$ et $\lambda = 3\,\text{cm}$~:
$\sigma_{2D}/\lambda \approx 2{,}2$, ce qui sert de validation du code.

% ══════════════════════════════════════════════════════════════════════════════
\section{Couche absorbante --- CPML}
% ══════════════════════════════════════════════════════════════════════════════

\subsection{Principe de la PML de Bérenger}

Pour simuler un milieu infini, on entoure le domaine de couches
\emph{Perfectly Matched Layers} (PML). La PML de Bérenger (1994) est
un milieu artificiel anisotrope à conductivité complexe, dont l'impédance
est identique à celle du vide quel que soit l'angle d'incidence~:
\begin{equation}
    \tilde{n}_{\text{PML}}(\omega) = n_0
    \quad \text{pour tout } \omega \text{ et tout } \theta
\end{equation}

\subsection{Formulation CPML}

La \emph{Convolutional PML} (CPML, Roden \& Gedney 2000) étend la PML à
l'aide d'un étirement de coordonnées complexe~:
\begin{equation}
    \tilde{\partial}_x \to \frac{1}{\kappa_x} \partial_x,
    \qquad \kappa_x = \kappa_x^{\max}\,, \quad \kappa_x^{\max} \in [1,\infty)
\end{equation}
et d'une convolution $\psi$ qui implémente le terme de mémoire~:

\begin{align}
    \psi^{n+1} &= b\,\psi^n + a\,\frac{\partial E_z}{\partial x}
    \label{eq:psi_update} \\
    b &= \exp\!\left[-\!\left(\frac{\sigma}{\kappa}+\alpha\right)\frac{\dt}{\eps_0}\right] \\
    a &= \frac{\sigma}{\sigma\kappa + \kappa^2\alpha}\,(b-1)
\end{align}

Le profil de conductivité suit une loi polynomiale d'ordre $m=3$~:
\begin{equation}
    \sigma(d) = \sigma_{\max}\left(\frac{d}{d_{\text{PML}}}\right)^m,
    \qquad
    \sigma_{\max} = \frac{0{,}8(m+1)}{\eta_0\,\dx}
    \label{eq:sigma_pml}
\end{equation}
où $d$ est la profondeur dans la PML.

L'équation de mise à jour de $\Ez$ en présence de la CPML devient~:
\begin{equation}
    {E_z}^{n+1}_{i,j} = {E_z}^n_{i,j} + C_{Ezh}
    \left[
        \frac{1}{\kappa_x}\,\frac{\partial \Hy}{\partial x} + \psi_{Hy}^{n+\half}
        - \frac{1}{\kappa_y}\,\frac{\partial \Hx}{\partial y} - \psi_{Hx}^{n+\half}
    \right]
    \label{eq:update_Ez_CPML}
\end{equation}

% ══════════════════════════════════════════════════════════════════════════════
\section{Transformation champ proche -- champ lointain (NTFF)}
% ══════════════════════════════════════════════════════════════════════════════

\subsection{DFT running}

Pendant la simulation, on accumule la Transformée de Fourier Discrète (DFT)
des champs sur une surface fermée $\Gamma$ entourant le diffuseur~:
\begin{align}
    \tilde{E}_z(\omega) &= \sum_{n=0}^{N} E_z(n\dt)\,e^{-j\omega n\dt}\,\dt
    \label{eq:DFT_Ez}\\
    \tilde{H}_x(\omega) &= \sum_{n=0}^{N} H_x\!\left[(n+\half)\dt\right]
                e^{-j\omega(n+\half)\dt}\,\dt
\end{align}

On accumule séparément les champs \emph{totaux} et les champs \emph{incidents},
puis on soustrait pour obtenir les champs \emph{diffusés}~:
$\tilde{E}_z^{\text{diff}} = \tilde{E}_z^{\text{tot}} - \tilde{E}_z^{\text{inc}}$.

\subsection{Théorème de réciprocité et potentiels de Huygens}

Le champ lointain diffusé dans la direction $\phi$ est calculé via les
potentiels équivalents de Huygens~:
\begin{align}
    \mathcal{N}_z(\phi) &=
    \oint_\Gamma \tilde{H}_y^{\text{diff}}\,e^{jk\bm{\hat{r}}\cdot\bm{r}'}\,dl'
    - \oint_\Gamma \tilde{H}_x^{\text{diff}}\,e^{jk\bm{\hat{r}}\cdot\bm{r}'}\,dl'
    \label{eq:Nz} \\
    \mathcal{L}_\phi(\phi) &=
    \oint_\Gamma \tilde{E}_z^{\text{diff}}
    \left(-\hat{n}_x\sin\phi + \hat{n}_y\cos\phi\right)
    e^{jk\bm{\hat{r}}\cdot\bm{r}'}\,dl'
    \label{eq:Lphi}
\end{align}

Le champ électrique lointain s'obtient alors~:
\begin{equation}
    E_z^{\text{far}}(\phi) = \omega\mu_0\,\mathcal{N}_z + k\,\mathcal{L}_\phi
    \label{eq:Ez_far}
\end{equation}

La SER bistatique normalisée vaut finalement~:
\begin{equation}
    \frac{\sigma_{2D}(\phi)}{\lambda} = \frac{k}{4\pi}
    \frac{\left|E_z^{\text{far}}(\phi)\right|^2}
         {\left|\tilde{E}_z^{\text{inc}}\right|^2}
    \label{eq:bistatic_SER}
\end{equation}

\subsection{Métrique de fitness}

Pour l'optimisation, on utilise une métrique simplifiée mais numériquement robuste~:
l'énergie du champ diffusé dans la zone SF en amont du mur~:
\begin{equation}
    \mathcal{F}(\bm{p}) = \sum_{(i,j)\in\Omega_{\text{SF}}} {E_z}^2(i,j,t_f)
    \label{eq:fitness}
\end{equation}
où $\bm{p} \in [-1,1]^{N_s}$ est le vecteur de $N_s$ paramètres géométriques
du mur et $t_f = N\dt$ est l'instant final de la simulation.

% ══════════════════════════════════════════════════════════════════════════════
\section{Paramétrisation du mur}
% ══════════════════════════════════════════════════════════════════════════════

\subsection{Représentation par segments de contrôle}

Le profil du mur est défini par $N_s$ points de contrôle uniformément
répartis sur sa hauteur $H$. Le déplacement de la surface en chaque point
$j$ de la grille est calculé par interpolation linéaire~:
\begin{equation}
    d(j) = \Delta_{\max} \cdot
    \text{interp}\!\left(\frac{j-j_{\min}}{j_{\max}-j_{\min}},
    \; [0,\dots,1], \; [p_1,\dots,p_{N_s}]\right)
    \label{eq:wall_profile}
\end{equation}
avec $\Delta_{\max} = 10\,\text{cellules}$ le déplacement maximal autorisé.

\begin{figure}[h!]
\centering
\begin{tikzpicture}[scale=0.8]
    % Mur plat de référence
    \draw[gray!50, dashed, thick] (3,-3) -- (3,3) node[above,gray]{\small Plat};
    
    % Définition des points pour relier exactement les points de contrôle
    \coordinate (P0) at (3, -3);
    \coordinate (P1) at ({3 + 1.2*(-0.3)}, -2.5);
    \coordinate (P2) at ({3 + 1.2*(0.8)}, -1.5);
    \coordinate (P3) at ({3 + 1.2*(-0.9)}, -0.5);
    \coordinate (P4) at ({3 + 1.2*(1.0)}, 0.5);
    \coordinate (P5) at ({3 + 1.2*(-0.5)}, 1.5);
    \coordinate (P6) at ({3 + 1.2*(0.6)}, 2.5);
    \coordinate (P7) at (3, 3);
    
    % Remplissage (parfaitement ajusté aux points)
    \fill[blue!20, opacity=0.5]
        (4.5,-3) -- (P0) -- (P1) -- (P2) -- (P3) -- (P4) -- (P5) -- (P6) -- (P7) -- (4.5,3) -- cycle;
        
    % Mur optimisé (profil interpolé linéairement)
    \draw[blue!70!black, thick] (P0) -- (P1) -- (P2) -- (P3) -- (P4) -- (P5) -- (P6) -- (P7);
        
    % Points de contrôle
    \foreach \y/\p/\idx in {-2.5/-0.3/1, -1.5/0.8/2, -0.5/-0.9/3, 0.5/1.0/4, 1.5/-0.5/5, 2.5/0.6/6}
        \fill[red] ({3 + 1.2*\p}, \y) circle (3pt)
            node[right, font=\tiny, black, xshift=2pt]{$p_{\idx}$};
            
    % Onde incidente
    \foreach \yy in {-2,0,2}
        \draw[->, blue, thick, dashed] (0,\yy) -- (2.5,\yy);
    \node[blue] at (1.2, 2.5) {\small $\vec{k}_{\text{inc}}$};
    
    % Axe
    \draw[->] (-0.5,-3.5) -- (-0.5,3.5) node[left]{\small $y$};
    \draw[->] (-0.5,-3.5) -- (5.5,-3.5) node[below]{\small $x$};
    
    % Annotations
    \draw[gray!50, densely dotted, thick] (4.2, 0.5) -- (4.2, -3.3); % Ligne d'alignement pour le max (p_4)
    \draw[<->, thick] (3,-3.3) -- (4.2,-3.3) node[midway,below]{\tiny $\Delta_{\max}$};
    \node[blue!70!black, anchor=west] at (4.3,-3.3){\small Mur PEC};
\end{tikzpicture}
\caption{Paramétrisation du mur par $N_s = 6$ points de contrôle $p_i \in [-1,1]$.
Le profil est interpolé linéairement entre les points. L'onde plane incidente
arrive en incidence normale depuis la gauche.}
\label{fig:wall_parametrization}
\end{figure}

% ══════════════════════════════════════════════════════════════════════════════
\section{Algorithme Génétique}
% ══════════════════════════════════════════════════════════════════════════════

\subsection{Formulation générale}

L'algorithme génétique cherche~:
\begin{equation}
    \bm{p}^* = \arg\min_{\bm{p} \in [-1,1]^{N_s}} \mathcal{F}(\bm{p})
    \label{eq:opt_problem}
\end{equation}
Il maintient une population de $N_{\text{pop}}$ individus
$\{\bm{p}^{(k)}\}_{k=1}^{N_{\text{pop}}}$ et les fait évoluer sur $N_g$ générations.

\subsection{Initialisation par Latin Hypercube Sampling}

L'initialisation par LHS (McKay \textit{et al.}, 1979) garantit une
couverture uniforme de l'hypercube $[-1,1]^{N_s}$ supérieure à
l'échantillonnage aléatoire naïf.

Pour $N_{\text{pop}}$ individus et $N_s$ dimensions, on découpe chaque
dimension en $N_{\text{pop}}$ strates équiprobables de largeur $1/N_{\text{pop}}$
et on tire exactement un point par strate~:
\begin{equation}
    p_{k,j} = \frac{\pi_j(k) + U_{k,j}}{N_{\text{pop}}}\,(b-a) + a
    \label{eq:LHS}
\end{equation}
où $\pi_j$ est une permutation aléatoire et $U_{k,j} \sim \mathcal{U}(0,1)$.

\subsection{Sélection par tournoi}

À chaque génération, deux parents sont sélectionnés indépendamment.
La sélection par tournoi de taille $k$ consiste à tirer $k$ candidats
uniformément au hasard et à retenir le meilleur~:
\begin{equation}
    \bm{p}^{\text{parent}} = \arg\min\left\{\mathcal{F}(\bm{p}^{(i_1)}),
    \dots, \mathcal{F}(\bm{p}^{(i_k)})\right\}
    \label{eq:tournament}
\end{equation}
La pression sélective croît avec $k$~; on utilise $k=3$.

\subsection{Croisement SBX}

Le \emph{Simulated Binary Crossover} (SBX, Deb \& Agrawal 1995) produit deux
enfants $o_1, o_2$ à partir de deux parents $x_1 \leq x_2$ en suivant une
distribution polynomiale bornée dont la forme imite le croisement mono-point binaire.

Pour chaque gène, avec probabilité $1/2$, on calcule~:
\begin{equation}
    \beta_q = \begin{cases}
        (u\alpha)^{1/({\eta_c+1})} & \text{si } u \leq 1/\alpha \\
        (2-u\alpha)^{-1/({\eta_c+1})} & \text{sinon}
    \end{cases}
    \label{eq:sbx_beta}
\end{equation}
avec $u \sim \mathcal{U}(0,1)$,
$\alpha = 2 - \beta^{-(\eta_c+1)}$,
$\beta = 1 + 2(x_1 - l)/(x_2 - x_1)$,
puis~:
\begin{align}
    o_1 &= \tfrac{1}{2}\left[(x_1+x_2) - \beta_q(x_2-x_1)\right] \\
    o_2 &= \tfrac{1}{2}\left[(x_1+x_2) + \beta_q(x_2-x_1)\right]
    \label{eq:sbx_children}
\end{align}

L'indice $\eta_c = 10$ contrôle l'étalement~: plus $\eta_c$ est grand,
plus les enfants restent proches des parents (exploitation vs exploration).

\subsection{Mutation polynomiale}

La mutation polynomiale bornée (Deb \& Goyal 1996) applique à chaque gène
avec probabilité $p_m = 1/N_s$ une perturbation~:
\begin{equation}
    y_i = x_i + \delta_i \cdot (b-a)
    \label{eq:pm_mutation}
\end{equation}
où $\delta_i$ suit la distribution~:
\begin{equation}
    \delta_i = \begin{cases}
        \left[2u + (1-2u)(1-\delta_L)^{\eta_m+1}\right]^{1/(\eta_m+1)} - 1
        & u < 0{,}5 \\
        1 - \left[2(1-u) + 2(u-0{,}5)(1-\delta_R)^{\eta_m+1}\right]^{1/(\eta_m+1)}
        & u \geq 0{,}5
    \end{cases}
    \label{eq:pm_delta}
\end{equation}
avec $\delta_L = (x_i - a)/(b-a)$, $\delta_R = (b-x_i)/(b-a)$.

\subsection{Règle du 1/5 de Rechenberg --- adaptation de $\eta_m$}

La règle du 1/5 de Rechenberg (1973) postule que le taux de succès optimal
des mutations est $1/5$~: si le taux observé sur une fenêtre glissante de $W$
mutations dépasse $1/5$, les mutations sont trop conservatives et on doit
augmenter l'exploration~; sinon on les réduit~:
\begin{equation}
    \eta_m \leftarrow \begin{cases}
        \max(\eta_m^{\min},\; \eta_m \cdot c) & \text{si } r_{\text{succ}} > 1/5 \\
        \min(\eta_m^{\max},\; \eta_m / c) & \text{si } r_{\text{succ}} \leq 1/5
    \end{cases}
    \label{eq:rechenberg}
\end{equation}
avec $c = 0{,}85$ et les bornes $\eta_m \in [2, 200]$.
Un $\eta_m$ élevé correspond à des mutations de faible amplitude (exploitation) ;
un $\eta_m$ faible à de fortes perturbations (exploration).

\subsection{Élitisme et Hall of Fame}

L'élitisme garantit que les $N_e$ meilleurs individus de la génération $g$
sont transmis sans modification à la génération $g+1$. Le \emph{Hall of Fame}
maintient en parallèle une archive des $N_{\text{HOF}}$ meilleurs individus
\emph{jamais} évalués (sur toutes les générations), indépendamment de la
dynamique de la population.

\subsection{Détection de stagnation et redémarrage adaptatif}

La diversité de la population est mesurée par l'écart-type moyen normalisé
des génomes~:
\begin{equation}
    D = \frac{1}{N_s} \sum_{j=1}^{N_s}
    \frac{\text{std}_{k}\!\left[p_{k,j}\right]}{(b-a)/2}
    \in [0,1]
    \label{eq:diversity}
\end{equation}

Un redémarrage est déclenché si~: (i) $D < D_{\min} = 0{,}02$ (convergence prématurée),
ou (ii) $\mathcal{F}^* (g) - \mathcal{F}^*(g-W) < \eps_s$ (stagnation sur $W$
générations). Dans ce cas, $40\,\%$ de la population est remplacée par un
nouveau LHS, les $60\,\%$ meilleurs étant conservés.

\begin{algorithm}[H]
\caption{Algorithme Génétique --- boucle principale}
\label{alg:GA}
\begin{algorithmic}[1]
\Input $N_{\text{pop}}, N_g, N_s, \eta_c, \eta_m^0, p_m, p_c$
\Output Meilleur génome $\bm{p}^*$ et fitness $\mathcal{F}^*$
\State $\mathcal{P}_0 \leftarrow \text{LHS}(N_{\text{pop}}, N_s)$ \Comment{initialisation}
\State Évaluer $\mathcal{F}(\bm{p})$ pour tout $\bm{p} \in \mathcal{P}_0$
\State HOF $\leftarrow \{N_{\text{HOF}}\text{ meilleurs}\}$
\For{$g = 1, \dots, N_g$}
    \State $\mathcal{P}' \leftarrow \{N_e\text{ meilleurs de }\mathcal{P}_{g-1}\}$ \Comment{élites}
    \While{$|\mathcal{P}'| < N_{\text{pop}}$}
        \State $\bm{p}_1 \leftarrow \text{Tournoi}(\mathcal{P}_{g-1})$ ;
               $\bm{p}_2 \leftarrow \text{Tournoi}(\mathcal{P}_{g-1})$
        \If{$u \sim \mathcal{U}(0,1) < p_c$}
            \State $(o_1, o_2) \leftarrow \text{SBX}(\bm{p}_1, \bm{p}_2, \eta_c)$
        \Else
            \State $o_1 \leftarrow \bm{p}_1$ ; $o_2 \leftarrow \bm{p}_2$
        \EndIf
        \State $o_1 \leftarrow \text{PM}(o_1, \eta_m, p_m)$ ;
               $o_2 \leftarrow \text{PM}(o_2, \eta_m, p_m)$
        \State $\mathcal{P}' \leftarrow \mathcal{P}' \cup \{o_1, o_2\}$
    \EndWhile
    \State Évaluer $\mathcal{F}$ pour les nouveaux individus (parallèle)
    \State Mettre à jour HOF
    \State Adapter $\eta_m$ via règle du 1/5
    \If{stagnation ou $D < D_{\min}$}
        \State Redémarrage partiel (LHS sur $40\,\%$ de la pop)
    \EndIf
    \State $\mathcal{P}_g \leftarrow \mathcal{P}'$
\EndFor
\State \Return HOF.meilleur
\end{algorithmic}
\end{algorithm}

% ══════════════════════════════════════════════════════════════════════════════
\section{Étude de convergence numérique}
% ══════════════════════════════════════════════════════════════════════════════

\subsection{Convergence en résolution spatiale}

On mesure l'énergie rétrodiffusée d'un mur plat en variant la résolution
$\text{ppw} \in \{8, 10, 12, 15, 20\}$ (points par longueur d'onde).
Le tableau~\ref{tab:grid_conv} et la figure~\ref{fig:grid_conv} montrent
la stabilisation de l'énergie à partir de $\text{ppw} = 12$.

\begin{table}[h!]
\centering
\caption{Convergence en résolution spatiale ($\sum {E_z}^2$ mur plat).}
\label{tab:grid_conv}
\begin{tabular}{ccccc}
\toprule
ppw & $\dx$ (mm) & Taille grille totale & $\mathcal{F}$ & Variation relative \\
\midrule
8  & 3,75 & $104\times104$ & $\sim 3{,}2\times10^{-2}$ & --- \\
10 & 3,00 & $124\times124$ & $\sim 2{,}8\times10^{-2}$ & $-12{,}5\,\%$ \\
12 & 2,50 & $148\times148$ & $\sim 2{,}6\times10^{-2}$ & $-7{,}1\,\%$ \\
15 & 2,00 & $174\times174$ & $\sim 2{,}5\times10^{-2}$ & $-3{,}8\,\%$ \\
20 & 1,50 & $224\times224$ & $\sim 2{,}45\times10^{-2}$ & $-2{,}0\,\%$ \\
\bottomrule
\end{tabular}
\end{table}

\begin{figure}[h!]
\centering
\begin{tikzpicture}
\begin{axis}[
    width=0.7\textwidth, height=5.5cm,
    xlabel={Résolution (pts/$\lambda$)},
    ylabel={$\mathcal{F}$ (u.a.)},
    title={Convergence en résolution spatiale --- mur plat PEC},
    grid=both, grid style={gray!30},
    xtick={8,10,12,15,20},
    mark size=3pt,
    legend pos=north east,
]
\addplot[blue, thick, mark=*, mark options={fill=blue}]
    coordinates {(8,0.032)(10,0.028)(12,0.026)(15,0.025)(20,0.0245)};
\addlegendentry{FDTD ($\sum {E_z}^2$)}
\addplot[red, dashed, thick]
    coordinates {(8,0.0245)(20,0.0245)};
\addlegendentry{Valeur de référence (ppw=20)}
\end{axis}
\end{tikzpicture}
\caption{L'énergie rétrodiffusée converge asymptotiquement avec la résolution.
À 15 pts/$\lambda$, l'erreur relative est inférieure à 4\,\%.}
\label{fig:grid_conv}
\end{figure}

\subsection{Convergence en épaisseur de PML}

La CPML est testée pour $N_{\text{PML}} \in \{6, 8, 10, 12, 16\}$ cellules.
Dès $N_{\text{PML}} = 10$, le coefficient de réflexion résiduel de la PML
est inférieur à $-60\,\text{dB}$, confirmant l'absence de pollution numérique
aux frontières.

% ══════════════════════════════════════════════════════════════════════════════
\section{Résultats et discussion}
% ══════════════════════════════════════════════════════════════════════════════

\subsection{Résultats de l'optimisation GA}

\begin{table}[h!]
\centering
\caption{Comparaison des méthodes d'optimisation.}
\label{tab:results}
\begin{tabular}{lccc}
\toprule
\textbf{Méthode} & $\mathcal{F}$ & Réduction & Simulations FDTD \\
\midrule
Mur plat (baseline) & $\mathcal{F}_0$ & $0\,\%$ & 1 \\
Algorithme Génétique & $0{,}065\,\mathcal{F}_0$ & $\mathbf{93{,}5\,\%}$ & $\sim 5\,000$ \\
\bottomrule
\end{tabular}
\end{table}

\subsection{Nature du profil optimal}

La géométrie optimisée converge systématiquement vers un profil en \emph{dents
de scie} ou en \emph{zigzag}. Ce type de profil crée des \emph{cavités
résonantes} dont les dimensions sont de l'ordre de $\lambda/4$.
Le mécanisme physique est double~:
\begin{enumerate}
    \item \textbf{Interférences destructives~:} Les ondes partiellement
    réfléchies par les différentes facettes du profil se combinent en
    opposition de phase dans la direction rétrodiffusée ($\phi = \pi$).
    \item \textbf{Redirection angulaire~:} L'énergie qui n'est pas détruite
    par interférence est renvoyée dans des directions non dangereuses
    ($\phi \neq \pi$), hors du lobe principal du radar.
\end{enumerate}

\begin{figure}[h!]
\centering
\begin{tikzpicture}[scale=0.9]
    % Mur en dents de scie
    \filldraw[fill=gray!25, draw=black!70, thick]
        (3,-3) -- (2.2,-2) -- (3,-1) -- (2.2,0) -- (3,1) -- (2.2,2) -- (3,3)
        -- (4.5,3) -- (4.5,-3) -- cycle;
    % Ondes incidentes
    \foreach \yy in {-2.5,-1.5,-0.5,0.5,1.5,2.5}
        \draw[->, blue, thick] (0,\yy) -- (1.8,\yy);
    \node[blue,above] at (0.9,2.8){\small onde plane};
    % Flèches refléchies (dispersées)
    \draw[->, red!70, thick] (2.2,-2) -- (0.5,-3.0);
    \draw[->, red!70, thick] (2.2,-2) -- (0.5,-1.0);
    \draw[->, red!70, thick] (2.2,0)  -- (0.5, 1.0);
    \draw[->, red!70, thick] (2.2,0)  -- (0.5,-0.5);
    \draw[->, red!70, thick] (2.2,2)  -- (0.5, 3.0);
    \draw[->, red!70, thick] (2.2,2)  -- (0.5, 1.5);
    % Annotation
    \node[red!70, font=\small] at (-0.5,0){réflexions};
    \node[red!70, font=\small] at (-0.5,-0.5){dispersées};
    % Dimension lambda/4
    \draw[<->, green!50!black, thick] (4.7,1) -- (4.7,3)
        node[midway, right, font=\small]{$\sim\lambda/4$};
\end{tikzpicture}
\caption{Profil optimal en dents de scie obtenu par le GA. Les dents créent
des réflexions multiples dans des directions non rétrodiffusées. La hauteur de
chaque dent est de l'ordre de $\lambda/4$ pour maximiser les interférences
destructives en $\phi=\pi$.}
\label{fig:sawtooth}
\end{figure}

\subsection{Comparaison simulation/expérience}

Le rapport rapporte que le mur imprimé en 3D et
métallisé produit une atténuation confirmée mais inférieure aux prédictions
numériques. Deux sources d'écart ont été identifiées~:

\begin{enumerate}
    \item \textbf{Précision dimensionnelle~:} Le profil optimal est sensible
    à des tolérances inférieures à $\lambda/20 \approx 1{,}5\,\text{mm}$.
    Les imprimantes 3D standard ont des tolérances de $\pm 0{,}3\,\text{mm}$,
    mais les erreurs de métallisation peuvent atteindre $\pm 1\,\text{mm}$.

    \item \textbf{Hypothèse d'onde plane~:} La simulation injecte une onde
    plane \emph{infiniment large} en incidence normale. Le banc expérimental
    utilise deux cornets en champ proche, générant un front d'onde parabolique
    avec un angle d'ouverture non nul. Le code a été étendu pour gérer
    plusieurs angles d'incidence ($0°$, $\pm 15°$) afin d'améliorer la
    robustesse angulaire de la solution.
\end{enumerate}

\begin{remarque}{Remarque --- Validité du modèle 2D}
La simulation est strictement 2D (invariance selon $z$). En réalité, le mur
a une dimension finie selon $z$, ce qui produit des effets de diffraction
aux bords. Pour corriger cela, il faudrait une simulation FDTD 3D, dont le
coût est $O(N_x N_y N_z N_t)$ au lieu de $O(N_x N_y N_t)$~: environ $N_z \sim 50$
fois plus coûteux.
\end{remarque}

% ══════════════════════════════════════════════════════════════════════════════
\section{Conclusion}
% ══════════════════════════════════════════════════════════════════════════════

Ce projet démontre la faisabilité d'une optimisation de géométrie anti-radar
entièrement guidée par la physique électromagnétique simulée. Les points clés
sont~:

\begin{enumerate}
    \item \textbf{Physique rigoureuse~:} Le solveur FDTD TMz implémente
    fidèlement les équations de Maxwell avec CPML, TFSF et NTFF, des méthodes
    de référence en simulation électromagnétique.

    \item \textbf{Résultats significatifs~:} Des réductions de SER de l'ordre
    de $93$--$94\,\%$ sont obtenues de manière reproductible sur le mur plat
    de référence.

    \item \textbf{Lien expérimental~:} La validation sur banc hyperfréquences
    confirme qualitativement les prédictions, et quantifie les écarts dus aux
    limitations matérielles (tolérances 3D, champ proche, invariance 2D).
\end{enumerate}

\paragraph{Perspectives.}
Les extensions naturelles comprennent~: (i) simulation 3D pour lever
l'hypothèse d'invariance~; (ii) co-optimisation géométrie + RAM (matériaux
absorbants)~; (iii) robustesse angulaire étendue à un demi-espace frontal~;
(iv) intégration d'un réseau de neurones pour accélérer
l'évaluation de la fitness.


% ══════════════════════════════════════════════════════════════════════════════
\addcontentsline{toc}{section}{Références}

\begin{thebibliography}{99}

\bibitem{yee1966}
K.S. Yee, ``Numerical solution of initial boundary value problems involving
Maxwell's equations in isotropic media,''
\textit{IEEE Transactions on Antennas and Propagation}, vol.~14, no.~3,
pp.~302--307, 1966.

\bibitem{taflove2005}
A. Taflove et S.C. Hagness,
\textit{Computational Electrodynamics: The Finite-Difference Time-Domain Method},
3\textsuperscript{e} éd. Artech House, 2005.

\bibitem{berenger1994}
J.-P. Bérenger, ``A perfectly matched layer for the absorption of
electromagnetic waves,''
\textit{Journal of Computational Physics}, vol.~114, no.~2, pp.~185--200, 1994.

\bibitem{roden2000}
J.A. Roden et S.D. Gedney, ``Convolutional PML (CPML): An efficient FDTD
implementation of the CFS-PML for arbitrary media,''
\textit{Microwave and Optical Technology Letters}, vol.~27, no.~5,
pp.~334--339, 2000.

\bibitem{balanis2012}
C.A. Balanis,
\textit{Advanced Engineering Electromagnetics}, 2\textsuperscript{e} éd.
Wiley, 2012.

\bibitem{deb1995}
K. Deb et R.B. Agrawal, ``Simulated binary crossover for continuous search
space,'' \textit{Complex Systems}, vol.~9, no.~2, pp.~115--148, 1995.

\bibitem{deb1996}
K. Deb et M. Goyal, ``A combined genetic adaptive search (GeneAS) for
engineering design,'' \textit{Computer Science and Informatics}, vol.~26,
no.~4, pp.~30--45, 1996.

\bibitem{rechenberg1973}
I. Rechenberg,
\textit{Evolutionstrategie: Optimierung technischer Systeme nach Prinzipien
der biologischen Evolution}. Frommann-Holzboog, Stuttgart, 1973.

\bibitem{mckay1979}
M.D. McKay, R.J. Beckman et W.J. Conover, ``A comparison of three methods for
selecting values of input variables in the analysis of output from a computer
code,'' \textit{Technometrics}, vol.~21, no.~2, pp.~239--245, 1979.

\bibitem{baker1985}
J.E. Baker, ``Adaptive selection methods for genetic algorithms,''
\textit{Proc. 1st International Conference on Genetic Algorithms}, pp.~101--111,
1985.

\bibitem{hansen2001}
N. Hansen et A. Ostermeier, ``Completely derandomized self-adaptation in
evolution strategies,'' \textit{Evolutionary Computation}, vol.~9, no.~2,
pp.~159--195, 2001.

\bibitem{williams1992}
R.J. Williams, ``Simple statistical gradient-following algorithms for
connectionist reinforcement learning,'' \textit{Machine Learning}, vol.~8,
no.~3--4, pp.~229--256, 1992.

\end{thebibliography}

\end{document}
